% Bachelorarbeit: Auswahl geeigneter 6R-Modernisierungsstrategien für frag.jetzt
% Sprache: formell und wissenschaftlich

\setlength{\emergencystretch}{2em}

\section{Einleitung}

Softwareprojekte stehen heute nicht mehr nur vor der Herausforderung, funktional zu sein, sondern auch langfristig effizient, nachhaltig und anpassungsfähig zu bleiben.
Als Lösung dafür wird die Migration von Client-Side-Software in die Cloud gesehen, welche wichtige Vorteile wie Skalierbarkeit und Flexibilität bietet, aber gleichzeitig
völlig neue Herausforderungen und Probleme mit sich bringt. Diesen Migrationsprozess zu gestalten und zu steuern, ist sehr Komplex. Weswegen damit fundierte Entscheidungen 
getroffen werden können, fundierte Fachkenntnisse und vorallem die Richtige Strategie zur Migration in die Cloud benötigt wird. Einher mit einer Migration geht meist eine Modernisierung der Software.
Die Frage für Organisationen ist nicht mehr \emph{ob}, sondern \emph{wie} bestehende Software modernisiert werden soll um eine möglichst schnelle,
kostengünstige aber auch effektive Cloud-Migration durchzuführen.

Verbreitete Modelle zur Cloud Migration werden in den sogenannten \enquote{6R}-Strategien beschrieben, welche jeweiles eine etwas ander Herangehensweise beschreiben.
Ziel dieser Arbeit ist eine begründete Entscheidung \emph{pro Komponente} zu liefern. Denn je nach technischem Zustand, Nutzen, Risiko und Kopplungsgrad soll abgewogen werden, 
welche der Strategien für die einzelnen Komponenten des Praxisbeispiels frag.jetzt am besten geeigent ist. Diese zuvor als \enquote{6R}-Strategien bezeichneten Vorgehensweisen 
beinhalten \emph{Rehost} (Lift-and-Shift), \emph{Replatform}, \emph{Refactor} bzw. \emph{Re-architect}, \emph{Repurchase/Replace}, \emph{Retire} oder \emph{Retain} ~\cite{aws2023sixRs}.

Diese Bachelorarbeit untersucht anhand der Anwendung \texttt{frag.jetzt}, welche der \enquote{6R}-Strategien für die jeweiligen Komponenten am besten geeignet ist. 
Die Software frag.jetzt lässt sich in klar unterscheidbare Komponenten zerlegen. Das Frontend wird als Progressive Web App (PWA) abgebildet, 
während das Backend die Anwendungslogik und API-Schicht bildet und dazu Java Spring-Boot verwendet. 
Zur Datenhaltung kommt eine PostgreSQL-Datenbank zu Einsatz. Ergänzt werden die Komponenten durch die angebundenen AI-Services welche KI-Funktionalitäten für die Nutzer bereitstellen.

\subsection{Forschungsfrage und Zielsetzung}
% bis hier überprüft
Die Arbeit beantwortet die Forschungsfrage:
\begin{quote}
\emph{Welche der 6~R-Strategien ist für die verschiedenen Komponenten von \texttt{frag.jetzt} jeweils am besten geeignet?}
\end{quote}

Ziel ist eine nachvollziehbar begründete Zuordnung einer Strategie \emph{pro Komponente}. Die Begründung erfolgt anhand eines Kriterienkatalogs, der sowohl technische 
als auch organisatorische Aspekte berücksichtigt, insbesondere Wartbarkeit, Sicherheitsanforderungen, Integrationsaufwand, Performance, Betriebskosten und erwartete Weiterentwicklung. 
Die Ergebnisse sollen so formuliert sein, dass sie als Entscheidungsgrundlage für eine schrittweise Modernisierung nutzbar sind.

\subsection{Aufbau der Arbeit}

Die Arbeit gliedert sich in vier Teile: (1) theoretische Fundierung der Modernisierungsstrategien und relevanter Qualitätsattribute, (2) Beschreibung des Systemkontexts und 
der betrachteten Komponenten, (3) methodisches Vorgehen zur Bewertung und Auswahl der Strategien sowie (4) Ergebnisdarstellung als begründete Zuordnung je Komponente inklusive 
Umsetzungsfahrplan und Reflexion der Grenzen.

\section{Theoretische Fundierung}

\subsection{Software-Modernisierung als architektonische Entscheidungsaufgabe}

Software-Modernisierung umfasst Maßnahmen, die technische und strukturelle Eigenschaften eines Systems verbessern, um langfristige Wartbarkeit, Anpassbarkeit und 
Betriebsfähigkeit zu sichern. Aus architektonischer Sicht ist Modernisierung eine Abfolge von Entscheidungen unter Unsicherheit: Jede Maßnahme verändert Kopplungen, 
Schnittstellen und Qualitätsmerkmale und kann dadurch neue Risiken oder Abhängigkeiten erzeugen. Architektur wird in diesem Sinne als \emph{Trade-off} zwischen konkurrierenden 
Qualitätszielen verstanden, etwa zwischen schneller Lieferung und hoher Robustheit oder zwischen Kostenminimierung und Reduktion von Vendor-Lock-in~\cite{richards2020fundamentals,ford2021hardparts}.

Ein zentraler Grundsatz ist, dass nicht alle Komponenten denselben Modernisierungsgrad benötigen. Komponenten unterscheiden sich bezüglich Änderungsrate, 
Sicherheitskritikalität, Skalierungsbedarf und Technologierisiko. Eine PWA ist beispielsweise stark durch Browser-APIs, Endgerätevielfalt und UX-Anforderungen geprägt, während eine Datenbank 
primär durch Datenintegrität, Migrationsrisiken und Betriebsstabilität dominiert wird. Folglich ist eine komponentenspezifische Strategieauswahl plausibler als
 eine monolithische \enquote{One-size-fits-all}-Modernisierung.

\subsection{Die 6~R-Strategien}

Das 6~R-Modell dient der strukturierten Portfolio-Entscheidung. AWS beschreibt einen eng verwandten Katalog von Migrationsstrategien (häufig als \enquote{7~Rs} bezeichnet);
 in dieser Arbeit wird für die weitere Argumentation eine gebräuchliche 6~R-Variante verwendet, in der die zusätzliche Kategorie \enquote{Relocate/Verschieben} nicht separat betrachtet wird~\cite{aws2023sixRs}:
\begin{itemize}
  \item \textbf{Rehost:} Migration mit minimalen Änderungen (z.\,B. Infrastrukturwechsel), primär zur schnellen Entlastung oder Standardisierung.
  \item \textbf{Replatform:} Migration mit begrenzten technischen Anpassungen, typischerweise zur Nutzung eines managed Dienstes oder zur Reduktion operativer Last.
  \item \textbf{Refactor / Re-architect:} Strukturelle Änderung der Anwendung (z.\,B. Modularisierung, Entkopplung, Cloud-Native-Patterns) zur Verbesserung von Skalierbarkeit, Wartbarkeit oder Resilienz.
  \item \textbf{Repurchase / Replace:} Wechsel zu einem anderen Produkt (z.\,B. SaaS) statt Weiterbetrieb der bestehenden Lösung.
  \item \textbf{Retire:} Außerbetriebnahme, wenn kein geschäftlicher Nutzen (mehr) besteht.
  \item \textbf{Retain:} Beibehaltung (vorerst keine Migration), z.\,B. bei Abhängigkeiten oder fehlender Priorität.
\end{itemize}

In der Praxis werden einzelne Begriffe unterschiedlich abgegrenzt (insbesondere \emph{Refactor} vs. \emph{Re-architect}). In dieser Arbeit werden beide als Spektrum
 tiefgreifender, architektureller Veränderungen verstanden, die über rein kosmetische Code-Verbesserungen hinausgehen.

\subsection{Komponententypen und Qualitätsanforderungen}

Für die Bewertung der Strategien werden Qualitätsanforderungen herangezogen, die in der Softwarearchitektur als entscheidungsleitend gelten~\cite{richards2020fundamentals,ford2021hardparts}. 
Für \texttt{frag.jetzt} sind insbesondere relevant:
\begin{itemize}
  \item \textbf{Wartbarkeit und Änderbarkeit:} Aufwand zur Implementierung neuer Anforderungen ohne regressionsbedingte Risiken.
  \item \textbf{Sicherheit und Datenschutz:} Schutz vor Missbrauch, Einhaltung organisatorischer und rechtlicher Vorgaben.
  \item \textbf{Betriebsstabilität (Resilienz):} Fähigkeit, Fehler zu tolerieren, degradierte Zustände zu beherrschen und Wiederherstellung zu ermöglichen.
  \item \textbf{Kosten und Effizienz:} direkte Betriebskosten sowie indirekte Kosten durch Komplexität und organisatorischen Overhead.
  \item \textbf{Portabilität und Lock-in-Risiko:} Abhängigkeit von spezifischen Plattformen oder Diensten.
\end{itemize}

Im Kontext von AI-Services treten zusätzliche Risiken auf, etwa durch Modell-Updates, unvorhersehbare Antwortqualität, Datenschutzanforderungen und die Notwendigkeit 
systematischer Risikosteuerung. Rahmenwerke wie der NIST AI RMF bieten hierfür Kategorien zur Risikoidentifikation und -behandlung, die sich als ergänzende Kriterien operationalisieren lassen~\cite{nist2023airmf}.

\subsection{Systemkontext: frag.jetzt als Fallstudie}

\texttt{frag.jetzt} wird in dieser Arbeit als exemplarisches System verstanden, an dem sich die Strategieauswahl für typische Webanwendungs-Komponenten 
demonstrieren lässt. Die technische Aufteilung in Frontend (PWA), Backend (Spring Boot), Datenhaltung (PostgreSQL) und externe AI-Services ist in modernen Anwendungen häufig 
anzutreffen und ermöglicht eine getrennte Betrachtung. Für die Komponentenbeschreibung werden primär Architekturartefakte (z.\,B. Schnittstellen, Deployments, Datenflüsse) und 
betrieblich relevante Eigenschaften (z.\,B. Skalierungsprofile, Ausfallfolgen) herangezogen.

\section{Methodik}

\subsection{Forschungsdesign}

Methodisch wird eine qualitative Fallstudie mit einem formalisierten Bewertungsrahmen kombiniert. Die Fallstudie erlaubt es, kontextabhängige 
Faktoren (z.\,B. organisatorische Constraints, Betriebsrealität, technische Historie) zu berücksichtigen. Der Bewertungsrahmen sorgt dafür, dass die resultierende Strategieentscheidung 
nicht rein intuitiv erfolgt, sondern entlang expliziter Kriterien transparent begründet wird.

\subsection{Bewertungsrahmen: Kriterienkatalog und Entscheidungsmatrix}

Für jede Komponente wird eine Entscheidungsmatrix erstellt, in der die 6~R-Strategien entlang eines Kriterienkatalogs bewertet werden.
 Kriterien werden in mess- oder zumindest begründbar beobachtbare Indikatoren überführt (z.\,B. \enquote{Deployment-Frequenz} als Proxy für Änderungsrate, \enquote{Anzahl externer Abhängigkeiten} als Proxy
  für Integrationskomplexität). Die Bewertung erfolgt in einer ordinalen Skala (z.\,B. 1--5), ergänzt durch qualitative Begründungen.

Ein Beispiel für die Struktur einer solchen Matrix zeigt Tabelle~\ref{tab:entscheidung}.

\begin{table}[ht]
\centering
\footnotesize
\resizebox{\columnwidth}{!}{%
\begin{tabular}{@{}lcccc@{}}
\toprule
\textbf{Kriterium} & \textbf{Rehost} & \textbf{Replatform} & \textbf{Refactor} & \textbf{Replace} \\
\midrule
Betriebsrisiko & mittel & mittel & hoch & niedrig--mittel \\
Kosten (kurzfristig) & niedrig & niedrig--mittel & hoch & mittel \\
Wartbarkeit (langfristig) & niedrig & mittel & hoch & mittel--hoch \\
Lock-in-Risiko & niedrig & mittel & mittel & hoch \\
\bottomrule
\end{tabular}
}
\caption{Schematische Entscheidungsmatrix (Beispielstruktur; Werte komponentenspezifisch)}
\label{tab:entscheidung}
\end{table}

Die Gewichtung der Kriterien wird explizit festgelegt (z.\,B. durch Experteneinschätzung oder vereinfachte Paarvergleiche). Dadurch kann transparent gezeigt werden, warum eine Strategie trotz hoher Implementierungskosten dennoch vorzugswürdig sein kann (etwa bei hohen Sicherheitsanforderungen).

\subsection{Datenerhebung}

Zur Datenerhebung werden folgende Quellen vorgesehen:
\begin{itemize}
  \item \textbf{Architektur- und Codeanalyse:} Sichtung von Modulgrenzen, Schnittstellen, Datenmodellen, Abhängigkeiten und Build-/Deployment-Prozessen.
  \item \textbf{Betriebsdaten (sofern verfügbar):} Logs, Monitoring, Ausfallberichte, Latenz-/Lastprofile und Kostenindikatoren.
  \item \textbf{Dokumente und Standards:} technische Dokumentation zu PWA~\cite{mdn2024pwa}, Spring Boot~\cite{springboot2024docs} und PostgreSQL~\cite{postgresql2023docs} sowie Risiko- und Governance-Rahmen für AI-Services~\cite{nist2023airmf}.
  \item \textbf{Stakeholder-Interviews (optional):} Erhebung nicht-technischer Anforderungen (z.\,B. Releasezyklen, Compliance, Teamkompetenzen).
\end{itemize}

\subsection{Auswertung und Validität}

Die Auswertung erfolgt komponentenweise: Für jede Komponente werden (1) Ist-Zustand und Qualitätsanforderungen beschrieben, (2) Strategien anhand der Matrix bewertet und (3) 
eine Empfehlung mit Begründung und Risiken formuliert. Zur Erhöhung der Validität wird eine \emph{Plausibilitätsprüfung} eingeplant: Die Empfehlungen werden gegen 
bekannte Architekturprinzipien sowie gegen typische Fallstricke moderner Architekturen gespiegelt (z.\,B. Überfragmentierung, operativer Overhead, unbeabsichtigte 
Lock-in-Effekte)~\cite{richards2020fundamentals,ford2021hardparts}.

Grenzen der Methodik liegen vor allem in der Kontextabhängigkeit: Ergebnisse sind nicht ohne Weiteres auf andere Systeme übertragbar. Genau darin liegt jedoch zugleich 
der Beitrag einer Fallstudie, nämlich die nachvollziehbare Argumentation im konkreten Kontext.

\section{Geplante Gliederung der Arbeit}

Die folgende Gliederung beschreibt den vorgesehenen Aufbau der Bachelorarbeit und dient als Arbeitsplan. Kapitelnummern und Umfang sind als Zielwerte zu verstehen.

\begin{enumerate}
  \item \textbf{Einleitung}
  \begin{itemize}
    \item Problemstellung und Motivation
    \item Forschungsfrage, Zielsetzung, Beitrag
    \item Vorgehensüberblick
  \end{itemize}

  \item \textbf{Theoretische Grundlagen}
  \begin{itemize}
    \item Software-Modernisierung und Architekturentscheidungen~\cite{richards2020fundamentals,ford2021hardparts}
    \item 6~R-Strategien als Portfolio-Ansatz~\cite{aws2023sixRs}
    \item Qualitätsattribute und Bewertungskriterien
  \end{itemize}

  \item \textbf{Systemkontext und Komponentenbeschreibung (frag.jetzt)}
  \begin{itemize}
    \item Frontend-PWA: Architektur, Constraints und Release-/UX-Anforderungen~\cite{mdn2024pwa}
    \item Spring-Boot-Backend: API-Schicht, Domänenlogik, Integrationen~\cite{springboot2024docs}
    \item PostgreSQL: Datenmodell, Migrations- und Betriebsaspekte~\cite{postgresql2023docs}
    \item AI-Services: Abhängigkeiten, Datenschutz, Kosten und Risikoaspekte~\cite{nist2023airmf}
  \end{itemize}

  \item \textbf{Methodik}
  \begin{itemize}
    \item Kriterienkatalog, Gewichtung, Entscheidungsmatrix
    \item Datenerhebung und Auswertungsverfahren
    \item Validität, Limitationen
  \end{itemize}

  \item \textbf{Ergebnisse: Strategieempfehlung je Komponente}
  \begin{itemize}
    \item Begründete Zuordnung einer 6~R-Strategie pro Komponente
    \item Risikoanalyse und Annahmen
  \end{itemize}

  \item \textbf{Diskussion und Fazit}
  \begin{itemize}
    \item Implikationen für eine schrittweise Modernisierung
    \item Grenzen der Ergebnisse und Ausblick
  \end{itemize}
\end{enumerate}

\section*{Glossar}

\begingroup
\RaggedRight
\begin{description}[style=nextline]
  \item[6~R-Strategien:] Sammlung von Modernisierungsoptionen zur systematischen Entscheidung pro Anwendungskomponente (z.\,B. Rehost, Replatform, Refactor/Re-architect, Repurchase/Replace, Retire, Retain)~\cite{aws2023sixRs}.

  \item[AI-Service:] Externer oder interner Dienst, der KI-Funktionalität bereitstellt (z.\,B. LLM-API, Embedding-Service). Solche Dienste erfordern besonderes Risikomanagement, u.\,a. hinsichtlich Datenschutz, Nachvollziehbarkeit und Kosten~\cite{nist2023airmf}.

  \item[Backend (Spring Boot):] Serverseitige Anwendungskomponente, die Geschäftslogik und Schnittstellen (typisch REST/HTTP) bereitstellt. In dieser Arbeit repräsentiert das Backend die API- und Domänenschicht~\cite{springboot2024docs}.

  \item[Lock-in:] Bindung an spezifische Plattformen, APIs oder Dienste, die einen späteren Wechsel erschweren (technisch oder wirtschaftlich).

  \item[PWA (Progressive Web App):] Webanwendung, die nativeähnliche Eigenschaften (Offline-Fähigkeit, Installierbarkeit, Push) durch moderne Browser-Technologien erreicht~\cite{mdn2024pwa}.

  \item[PostgreSQL:] Relationales Datenbanksystem; im Kontext dieser Arbeit steht es für die persistente Datenhaltung inklusive Transaktions- und Migrationsaspekten~\cite{postgresql2023docs}.

  \item[Refactor/Re-architect:] Tiefgreifende Umgestaltung einer Komponente zur Verbesserung zentraler Qualitätsattribute, z.\,B. Modularisierung, Entkopplung oder Cloud-Native-Patterns~\cite{aws2023sixRs,ford2021hardparts}.

  \item[Rehost:] Migration mit minimalen Änderungen (Lift-and-Shift), primär zur schnellen Verlagerung und Standardisierung~\cite{aws2023sixRs}.

  \item[Replatform:] Migration mit begrenzten technischen Anpassungen, häufig um Betriebsaufwand zu reduzieren oder managed Dienste zu nutzen~\cite{aws2023sixRs}.

  \item[Repurchase/Replace:] Wechsel zu einem anderen Produkt (z.\,B. SaaS), wenn Eigenerstellung keinen strategischen Mehrwert liefert oder Betriebskosten unverhältnismäßig hoch sind~\cite{aws2023sixRs}.

  \item[Retain:] Beibehaltung einer Komponente (vorerst keine Migration), z.\,B. aus Sicherheitsgründen, wegen Abhängigkeiten oder fehlender Priorität~\cite{aws2023sixRs}.

  \item[Retire:] Außerbetriebnahme einer Komponente, wenn kein geschäftlicher Nutzen (mehr) besteht oder Betrieb unverhältnismäßig teuer ist~\cite{aws2023sixRs}.

  \item[Resilienz:] Fähigkeit eines Systems, Störungen zu tolerieren und sich zu erholen; häufig eng verbunden mit Observability, Redundanz und klaren Fehlergrenzen.

  \item[Technische Schulden:] Metapher für langfristige Folgekosten kurzfristiger Kompromisse (z.\,B. fehlende Tests, starke Kopplung, veraltete Abhängigkeiten), die Wartbarkeit und Änderbarkeit beeinträchtigen.

\end{description}
\endgroup

\section*{Quellen dieses Exposés}

\begingroup
\RaggedRight
\begin{itemize}
  \item \fullcite{aws2023sixRs}
  \item \fullcite{richards2020fundamentals}
  \item \fullcite{ford2021hardparts}
  \item \fullcite{davis2019cloudnativepatterns}
  \item \fullcite{mdn2024pwa}
  \item \fullcite{springboot2024docs}
  \item \fullcite{postgresql2023docs}
  \item \fullcite{nist2023airmf}
\end{itemize}
\endgroup
